% Compléments distincts du rapport ENSAE déposé ; actualisés le 2 novembre 2026.
% Chiffres : sorties_verif_inst/136.txt à 139.txt. Nom conservé pour les maîtres existants.
\clearpage
\part*{Compléments Institut après dépôt}
\addcontentsline{toc}{part}{Compléments Institut après dépôt}

\begin{cle}
Cette partie porte sur la version Institut et les analyses disponibles au 2 novembre.
Elle ne change pas les folios du rapport ENSAE cités dans le cours. Le rapport déposé décrit
$582$ pertes et $91$ excès ; les analyses Institut sur $2000$--$2026$ retiennent $570$ pertes et
$86$ excès. Toujours annoncer le périmètre et dire si le résultat est dans le rapport ou constitue
un complément. Le protocole de capital publié ne devient pas validé par l'ajout de ces diagnostics.
\end{cle}

\section{Descendre à l'entité : la variante retenue et ses limites}
\lequoi{Le registre observe les pertes déclarées, sans donner une exposition exhaustive. Une
faible relation entre actifs et incidents enregistrés ne prouve pas une faible relation entre
actifs et risque réel.}
\laformule{Pour une entité d'actifs $A$,
\[
\lambda(A)=\lambda_0\left(\frac{A}{A_0}\right)^{b_\lambda},\qquad
m(A)=\left(\frac{A}{A_m}\right)^b,\qquad
B(A)=\operatorname{VaR}_{99{,}5\%}\!\left(\sum_{n=1}^{N}m(A)\sum_{k\in S_n}Y_{n,k}\right).
\]
$A$ est le total des actifs, $A_0$ la moyenne géométrique des actifs du panel d'ajustement et
$\lambda_0$ la fréquence ajustée à cette taille. $A_m$ est la médiane des actifs associés aux
événements TIC du panel financier $2005$--$2022$, distinct des $570$ pertes cyber ;
$b$ règle l'échelle des pertes. $N$ suit une binomiale négative de moyenne $\lambda(A)$ et de
forme $r$ ; $S_n$ contient les piliers atteints et $Y_{n,k}$ leurs pertes élémentaires.
Le quantile est recalculé ; on ne divise pas directement le capital du périmètre financier.}
\cheznous{Le panel financier avec actifs connus porte sur $2005$--$2022$. La variante ajuste
la fréquence sur les firmes de plus de $100$~Md\$ : $6\,048$ firmes-années, $1\,402$ firmes,
$A_0=570$~Md\$, $b_\lambda=0{,}494$, erreur-type par firme $0{,}066$, et $r=0{,}594$.
À $20$~Md\$ d'actifs, $\lambda(A)=0{,}0209$/an et le facteur de sévérité vaut $0{,}854$.
L'élasticité de sévérité $b=0{,}087$ provient d'un ajustement tronqué sur les pertes TIC ;
son application aux pertes cyber-finance est une hypothèse de transfert. Les $484$ pertes
observées sous le seuil servent désormais de corps empirique, au lieu de pertes nulles.
Le métier n'est pas intégré au protocole retenu.}
\lepiege{Le calcul d'entité utilise la cascade indépendante sans choc commun P4. Le chiffre de
tête à quatre canaux utilise la marche avec ce choc ; les $32$ configurations par pilier utilisent
la marche sans ce choc. Une validation obtenue sur l'un de ces protocoles ne se transfère pas
automatiquement aux autres.}
\lepiege{L'estimation au-delà de $100$~Md\$ est extrapolée aux entités plus petites.
La restriction vise à réduire l'effet des fenêtres courtes sans démontrer une correction de
sélection : $57{,}8\,\%$ des grandes firmes ne sont observées qu'un an. La dispersion estimée
contient aussi l'hétérogénéité entre firmes.}
\laphrase{La variante retenue donne $24{,}1$~M\euro{} pour l'entité notionnelle non conforme,
contre $169$~M\euro{} dans la lecture historique du registre. Fréquence, dispersion et pertes
sous le seuil changent ensemble ; la baisse ne vient pas de l'exposant seul. C'est une
extrapolation conditionnelle, pas un capital individuel validé.}
\appui{Mémoire Institut, \S~9.5 ; scripts 136 et 137. Les anciennes lectures à exposants
$0{,}074$, $0{,}27$ et $1$ restent des comparaisons historiques ; $0{,}27$ n'est plus la lecture
retenue. Le rapport ENSAE déposé conserve ses chiffres historiques.}

\section{Lire une plage sans confondre les incertitudes}
\cheznous{Au point d'expert, les états de propagation C, PC et NC donnent respectivement
$13{,}2$, $17{,}7$ et $24{,}1$~M\euro{}, soit $+83\,\%$ de C à NC. La figure 9.3 est recalée
sur cette descente par le script 139. La plage aux $1\,024$ sommets est de $15{,}8$ à
$27{,}5$~M\euro{}. Le bootstrap des firmes, qui réestime intercept, pente et dispersion,
donne $8{,}3$ à $53{,}0$~M\euro{}, à sévérité et graine fixées.}
\lepiege{La première plage fait varier le sens des liens à force posée ; elle n'est ni un
intervalle de confiance ni une preuve d'extrema globaux. La seconde traite le calage de fréquence
sous le protocole du panel ; elle ne corrige pas la sélection du registre. On ne les fusionne
pas en un intervalle global. Les $6\,049$ à $20\,188$~M\euro{} concernent un autre protocole,
sur le périmètre financier représenté par le registre, non tout le secteur mondial.}
\laphrase{Une plage doit toujours dire ce qui varie et ce qui reste fixé. La nouvelle lecture
donne un ordre de grandeur, tout en montrant que la connaissance du registre domine la précision
d'un chiffre isolé.}
\appui{Mémoire Institut, \S~9.5--9.6 ; scripts 136, 137 et 139.}

\section{Capture des incidents : plusieurs risques pour le même registre}
\laformule{Avec une capture indépendante du montant et du nombre d'incidents,
\[
N_{\text{réel}}\sim\mathrm{NB}(\mu/q_{\mathrm{obs}},r),\qquad
N_{\text{enregistré}}\mid N_{\text{réel}}\sim\mathrm{Binom}(N_{\text{réel}},q_{\mathrm{obs}}),
\]
la loi observée est $\mathrm{NB}(\mu,r)$. Ici $\mu$ est la moyenne enregistrée et
$q_{\mathrm{obs}}$ la probabilité de capture. Le registre seul ne sépare pas ces deux quantités.}
\cheznous{Pour la notionnelle, les captures supposées de $100$, $50$, $25$ et $10\,\%$
donnent des VaR de $23{,}1$, $74{,}0$, $153{,}9$ et $327{,}7$~M\euro{}, moyennes sur quatre
graines. Les ES correspondantes valent $160{,}6$, $282{,}4$, $468{,}5$ et $850{,}9$~M\euro{}.
Les $23{,}1$~M\euro{} se comparent aux $24{,}1$ d'une graine, au bruit de simulation près.}
\lepiege{Ces captures sont des scénarios, jamais des taux estimés. Les pertes publiques ont
probablement une capture liée à leur montant : l'indépendance est une hypothèse forte.
Le périmètre financier est concerné lui aussi ; corriger les seules entités fausserait la comparaison.}
\laphrase{Les $24$ millions ne sont ni une borne prudente ni un risque réel identifié :
plusieurs fréquences réelles peuvent produire le même registre observé.}
\appui{Mémoire Institut, \S~9.5.2 et 9.6.2 ; script 138.}

\section{Un quantile nul n'est pas une absence d'incident}
\laformule{Pour une perte annuelle non négative $L$,
\[
\Prob(L=0)\ge0{,}995\quad\Longrightarrow\quad
\operatorname{VaR}_{99{,}5\%}(L)=0.
\]
La condition porte sur la \emph{perte}, pas sur le seul nombre $N$ d'incidents. Sous la lecture
retenue, toutes les pertes d'incident sont positives, corps inclus ; la masse nulle vient donc
des années sans incident. Pour la binomiale négative,
$\Prob(L>0)=1-(r/(r+\lambda))^r$. L'ancien moteur annulait les pertes sous le seuil : il pouvait,
en plus, compter un incident sans perte positive.}
\cheznous{La bascule en forme fermée vaut $1\,076$~M\euro{} d'actifs. Les simulations donnent
$1\,053$ à $1\,145$~M\euro{} sur quatre graines ; le bootstrap des firmes donne $100$ à
$3\,889$~M\euro{}. À $500$~M\euro{} d'actifs, la petite entité a une VaR nulle mais une ES
d'environ $23$~M\euro{}. L'assureur A n'a plus une VaR nulle : elle vaut environ $1{,}0$~M\euro{}.}
\laphrase{La VaR nulle reflète la masse en zéro de la charge simulée à ce niveau de probabilité.
Un quantile plus élevé, ou une autre mesure de perte, peut informer autrement le risque.
Cette bascule ne définit pas la validité statistique du modèle.}
\appui{Mémoire Institut, \S~9.6 ; script 137 pour la bascule et son incertitude, 136 pour les entités.}

\section{Expliquer comment les quatre paramètres sont choisis}
\cheznous{La fréquence de base est comptée dans le registre. Le passage non conforme utilise
$\lambda_{NC}=\lambda_C\sum_v\pi_v m_v$ : $\pi_v$ est la part observée d'un vecteur d'attaque,
$m_v$ le centre d'une fourchette posée à partir d'effets publiés sur des contrôles. Le mélange
vaut environ $2{,}4875$, d'où $21{,}56$ à $53{,}62$ incidents par an. La détection multiplie
le taux $86/570=0{,}1509$ par $0{,}85$ ou $1{,}20$, facteurs posés. Les gains
$0{,}45/0{,}68/0{,}90$, la structure $T$ et le choc commun $\omega=0{,}68$ en NC sont posés.}
\lepiege{Ni les fourchettes des multiplicateurs ni les contraintes d'admissibilité sur $W$
ne sont des intervalles de confiance. La co-occurrence entre piliers n'est pas estimée dans
l'application. $T$ reste fixé, mais $W=gT/s_{\max}$ varie avec $g$. Le partage de graines
n'annule pas le bruit de simulation.}
\laphrase{Le calcul est reproductible ; le passage entre états reste un scénario. Il faut
annoncer séparément les données, les hypothèses et les sensibilités.}
\appui{Beamer ENSAE, diapo 9 en deux étapes ; mémoire Institut, chapitres 4 à 6 et 9.}

\section{Apprentissage statistique : un gain à vérifier hors période d'ajustement}
\lidee{Le script 131 demande si les caractéristiques observées améliorent une simple tendance
temporelle. Il prédit le dépassement de $u=20{,}03$~M\euro{} parmi les pertes enregistrées.
Ce n'est ni une fréquence annuelle d'incidents ni une mesure de conformité.}
\cheznous{Évaluation en origine glissante de $2014$ à $2025$ : $280$ pertes, $60$ dépassements.
M0 est un taux constant ; M1 une tendance logistique ; M2 une logistique pénalisée enrichie ;
M3 un boosting sur l'année, le métier, la région, la catégorie et les actifs. Les hyperparamètres
sont fixés avant l'évaluation ; l'incertitude est calculée par bootstrap des douze années évaluées.
Le gain de log-score de M3 contre M1 vaut $+0{,}0193$, intervalle à $90\,\%$
$[-0{,}0371\,;+0{,}0571]$. Aucun gain positif résolu n'est établi dans les variantes testées.}
\lepiege{Pour $37\,\%$ des pertes, les actifs décrivent une information financière postérieure
à l'incident. Les masquer ne fait pas apparaître de gain résolu. La variante GPD à covariables
attribue une densité nulle à un excès observé : sa moyenne sur les $59$ autres excès omet cet
échec de support et ne peut pas être présentée comme son score complet.}
\laphrase{Le boosting a le meilleur score moyen, mais le gain n'est pas distingué de zéro avec
la précision disponible. Ce diagnostic ne prouve pas que les caractéristiques sont inutiles ;
il ne justifie pas de remplacer ici le modèle de référence.}
\appui{Mémoire Institut, \S~4.7.1 ; script 131 et sortie 131, après dépôt ENSAE.}

\section{Descriptions SAS : observer un récit et identifier un mécanisme}
\lidee{Le script 132 est un pilote de $50$ descriptions ; le script 133 agrège le codage de
$570$ descriptions, pilote inclus. Au niveau strict, le texte établit une défaillance ; au
niveau large, le codeur ajoute une défaillance impliquée. Un seul assistant IA a codé ces textes.}
\cheznous{Le corpus complet contient $65$ incidents multi-piliers au sens strict et $203$ au sens
large. La corrélation de rang avec la matrice d'expert est de $0{,}67$ et $0{,}75$ pour les
co-occurrences brutes, mais de $0{,}33$ et $0{,}12$ pour les coefficients $\phi$, qui tiennent
compte des marges. Une forte fréquence de P1 contribue donc à l'accord brut apparent.}
\lepiege{Ces codes ne sont pas des défaillances mesurées et n'estiment ni $S$ ni $A$.
Aucune direction n'est codée. Un second codage est préparé sur $100$ descriptions, $50$ par strate,
avec questionnaire en aveugle et clé séparée. Les réponses ne sont pas encore reçues : aucun
accord intercodeurs n'est annoncé. Le pilote et le corpus complet ne sont pas deux validations indépendantes.}
\laphrase{Les textes peuvent aider à préciser ce qu'il faudrait observer. Avant de s'en servir
pour estimer un lien, il faut vérifier que deux lecteurs indépendants y repèrent les mêmes faits,
puis relier ces faits aux variables du modèle.}
\appui{Scripts 132--133 ; protocole \path{docs/protocoles/validation_codage_sas.md}.
Les textes sous licence et la clé ne font pas partie des fichiers à publier.}
